\subsection{The boundary torus and the actual affine-cover quotient}
\label{sec:ghost-generalized-jacobian}

\begin{result}{C.3.6}{Proved}{The complete compact fibre at two}
\label{lem:ghost-elliptic-two-fibre}
For every original A.44 output and its ordinary local point
\(t_2\), the compact three-adic Kummer classes of the
\(\mathbb Q_2\)-points of \(Y\) above \(t_2\) fill
\(E'(\mathbb Q_2)^{\wedge3}=\langle T\rangle\).
Consequently the compact local condition at two imposes
no additional restriction on the coefficients \(a,b\)
of \(a\kappa(P_0)+b\kappa(T)\).
\end{result}
\begin{proof}
Every exponent in A.44 is divisible by three, so
\(t_2/3\) and \((1-t_2)/(-2)\) are cubes in
\(\mathbb Q_2^\times\). Choose their cube roots
\(u,v\); they give a point of \(C\) above \(t_2\).
The proof of Proposition~\ref{prop:ghost-elliptic-selmer}
showed that \(E(\mathbb Q_2)\) is uniquely
three-divisible and that
\(\phi:E'(\mathbb Q_2)\to E(\mathbb Q_2)\)
is surjective, with kernel \(\langle T\rangle\).
The finite three-primary subgroup of \(E'(\mathbb Q_2)\)
must therefore lie in that kernel, and is exactly
\(\langle T\rangle\). The remaining open pro-two
subgroup is uniquely three-divisible. Thus its
three-adic completion is exactly this cyclic group.
The chosen point of \(C\) has three lifts through
\(\phi\), all in \(Y(\mathbb Q_2)\), and these
lifts differ by \(T\). Their Kummer classes fill
the asserted completion.
\end{proof}

Put \(D=E'\setminus Y=D_0\sqcup D_1\sqcup D_\infty\),
where the three terms lie over \(0,1,\infty\).
They are finite etale \(\mathbb Q\)-schemes of
degree nine. Write
\[
 K_i=\Gamma(D_i,\mathcal O_{D_i}),\qquad
 \mathcal J=\operatorname{Jac}(E',D),\qquad
 \mathcal T=(\operatorname{Res}_{D/\mathbb Q}\mathbb G_m)
                         /\mathbb G_m.
                                                        \tag{25}
\]
The notation \(K_i\) permits an etale algebra; the
norm below is its product norm. The Abel map
\(j:Y\to\mathcal J\) is based at the rational
origin \(O\in Y(\mathbb Q)\), which lies over
\(t=3\).

\begin{result}{C.3.7}{Proved}{The exact generalized-Jacobian image}
\label{prop:ghost-genjac-image}
Let \(K=\pi_1(Y_{\overline{\mathbb Q}})\), viewed
inside \(\pi_1(U_{\overline{\mathbb Q}})\) through
the specified degree-twenty-seven cover. The quotient
\(\pi_1(Y)/\overline{[K,K]}\) is the split
arithmetic extension with kernel \(T\mathcal J\),
split by \(O\); its inclusion into
\(\pi_1(U)/\overline{[K,K]}\) retains the cover
map. The image of its global continuous sections in
the original completed unit coordinates contains
\[
           (3C^{27},-2D^{27})
       \qquad(C,D\in\widehat{\mathbb Q^\times}).       \tag{26}
\]

More precisely, at coefficient three its image consists
exactly of the pairs
\[
 \begin{split}
 A&=3\left(u_0^a\frac{N_{K_0/\mathbb Q}\alpha_0}
                           {N_{K_\infty/\mathbb Q}\alpha_\infty}
      \right)^3,\\
 B&=-2\left(v_0^a\frac{N_{K_1/\mathbb Q}\alpha_1}
                            {N_{K_\infty/\mathbb Q}\alpha_\infty}
      \right)^3,                                      \\
 &a\in\mathbb Z_3,\quad
 \alpha_i\in\widehat{K_i^\times}_3,\\
 &u_0=\frac{253}{655},\quad v_0=-\frac{488}{655}.
 \end{split}\tag{27}
\]
All norms and powers in (27) are the continuous maps
on the indicated multiplicative completions.
\end{result}
\begin{proof}
We first identify the algebraic and fundamental-group
maps. The rigidified Picard description gives the exact
sequence of semiabelian varieties over \(\mathbb Q\)
\[
       1\longrightarrow\mathcal T
        \longrightarrow\mathcal J
        \longrightarrow E'\longrightarrow0.          \tag{28}
\]
This is Proposition 1.6 and equation (1.4.1) of
\cite{GhostJordanRibetScholl2024}, applied to the smooth
proper curve \(E'\) and reduced divisor \(D\).
The generalized Abel map is the semiabelian Albanese
map: \cite[Proposition A.8(ii),(iii)]{GhostMochizukiAlbanese2012}
identifies its geometric torsion-free abelianization
with \(T\mathcal J\). Here the geometric fundamental
group is free profinite of rank
\(1+27(2-1)=28\); its abelianization is already
torsion-free. Thus this is the actual geometric
abelianization, not a further quotient. Functoriality
and the point \(O\) identify the arithmetic quotient
with \(T\mathcal J\rtimes G_{\mathbb Q}\). The
splitting used here is the arithmetic point section
of \(O\). The subgroup \(K\) is the kernel of the
characteristic class-two, exponent-three Heisenberg
quotient of the geometric free group. It is therefore
normal in that group and stable under its arithmetic
action; hence \(\overline{[K,K]}\)
is normal in \(\pi_1(U)\), as required for the
displayed downstairs quotient.

On the pointed curve the two functions \(u,v\)
satisfy
\[
\begin{gathered}
 u(O)=v(O)=1,\\
 \operatorname{div}(u)=D_0-D_\infty,\quad
 \operatorname{div}(v)=D_1-D_\infty,\\
 t/3=u^3,\qquad (1-t)/(-2)=v^3.
\end{gathered}                                         \tag{29}
\]
The compact isogeny \(\phi\) is etale; the orders
one in the first two divisors follow already on
\(3u^3-2v^3=1\), and are retained under its pullback.
The pointed functions induce homomorphisms
\(U_0,V_0:\mathcal J\to\mathbb G_m\). To fix the
torus convention explicitly, represent a geometric
torus element with boundary values \((z_d)\) by
\([\operatorname{div}(h)]_D\), where
\(h(d)=z_d\). Approximation gives \(h\), and
two such choices differ by a function equal to one
on \(D\), so give the same generalized divisor
class; constant boundary values give zero. Weil
reciprocity with disjoint divisors gives
\[
 U_0|_{\mathcal T}(z)=
       \frac{\prod_{d\in D_0}z_d}{\prod_{d\in D_\infty}z_d},
 \qquad
 V_0|_{\mathcal T}(z)=
       \frac{\prod_{d\in D_1}z_d}{\prod_{d\in D_\infty}z_d}.
                                                        \tag{30}
\]
This derivation fixes the sign of both characters.
The formulas are Galois equivariant and descend to
the product norms on the etale algebras \(K_i\).

There is a \(\mathbb Q\)-homomorphism
\[
\begin{gathered}
 \beta:\mathbb G_m^2\longrightarrow\mathcal T,\\
 (s,r)\longmapsto
     [s\text{ on }D_0,\ r\text{ on }D_1,
                             \ 1\text{ on }D_\infty].
\end{gathered}
                                                        \tag{31}
\]
Equations (29)--(31) prove the literal identities
\[
 (U_0,V_0)\beta=(s^9,r^9),\qquad
 (U_0^3,V_0^3)\beta=(s^{27},r^{27}).                 \tag{32}
\]
For completed classes \(C,D\), take their Kummer
cocycle in \(T\mathbb G_m^2\), apply
\(T\beta\), and include it in \(T\mathcal J\).
The crossed-homomorphism equation makes
\(g\mapsto (T\beta)(\kappa(C,D)(g))s_O(g)\)
a continuous section of the stated arithmetic
quotient. Its original coordinates are (26), since
the base section maps to \(3,-2\) and (32) gives
the changes. This proves (26) with every coefficient
prime retained.

For the exact three-adic image, first compute
\[
 H^1(G_{\mathbb Q},T_3\mathcal T)
   =\frac{\widehat{K_0^\times}_3\oplus
          \widehat{K_1^\times}_3\oplus
          \widehat{K_\infty^\times}_3}
          {\widehat{\mathbb Q^\times}_3},             \tag{33}
\]
where the denominator maps diagonally. Indeed the
Tate-module sequence of
\(\mathbb G_m\to\operatorname{Res}_D\mathbb G_m
\to\mathcal T\), followed by continuous Shapiro
and Kummer theory, gives its cohomology sequence.
The possible cokernel in (33) is the kernel of
restriction on \(H^2(G_{\mathbb Q},\mathbb Z_3(1))\)
to the etale algebra of \(D\). The finite Kummer
sequences give
\[
 H^2(G_{\mathbb Q},\mathbb Z_3(1))
               =T_3\operatorname{Br}(\mathbb Q).
\]
Here the finite \(H^1\) transitions
\(\mathbb Q^\times/(\mathbb Q^\times)^{3^{n+1}}
\to\mathbb Q^\times/(\mathbb Q^\times)^{3^n}\)
are surjective, so the derived-limit term is zero;
the Brauer transitions are multiplication by three.
The last Tate module is torsion-free. Restriction
followed by corestriction is multiplication by
\(\deg D=27\), so this restriction is injective.
This proves (33). The diagonal first map is also
injective, since restriction followed by norm on
\(\widehat{\mathbb Q^\times}_3\) is multiplication
by twenty-seven and this Kummer group is torsion-free.

Next, \(H^0(G_{\mathbb Q},T_3E')=0\), since
the rational three-primary torsion has order three.
Hence (28) gives an injection of (33) into
\(H^1(G_{\mathbb Q},T_3\mathcal J)\). Its map to
\(H^1(G_{\mathbb Q},T_3E')\) is surjective for a
concrete reason. Both \(P_0,T\) belong to
\(Y(\mathbb Q)\); the Kummer classes of
\(j(P_0),j(T)\in\mathcal J(\mathbb Q)\) map to
\(\kappa(P_0),\kappa(T)\), which generate that
whole continuous group by
Proposition~\ref{prop:ghost-elliptic-h1}.
Consequently every class in
\(H^1(G_{\mathbb Q},T_3\mathcal J)\) has the form
\[
\begin{gathered}
 a\kappa(j(P_0))+b\kappa(j(T))+\xi,\\
 a\in\mathbb Z_3,\quad b\in\{0,1,2\},\quad
 \xi\in H^1(G_{\mathbb Q},T_3\mathcal T).
\end{gathered}                                         \tag{34}
\]
The chosen representatives for \(b\) make the toric
term unique after the proper class has been fixed.

Substitution in the already verified cubic and
isogeny maps gives
\[
 \begin{split}
 &(u(P_0),v(P_0))=(253/655,-488/655),\\
 &t(P_0)=48582831/281011375,\qquad
 (u(T),v(T))=(1,1).
 \end{split}                                           \tag{35}
\]
The last equality follows from \(\phi(T)=O\) and
the pointed cubic origin \((1,1)\). Applying
(29)--(30) to (33)--(35) gives exactly (27).
Conversely every choice in (27) gives the cocycle
(34), hence an actual section of the quotient.
This proves equality of the stated images.
\end{proof}

In particular the ordinary factorial construction
with exponents \(27(j+4)!\) supplies infinite-support
pairs in (26). The generalized-Jacobian quotient of
this fixed cover therefore has global sections over
those pairs. This assertion concerns the explicitly
displayed quotient by \(\overline{[K,K]}\); lifting
its sections through that retained commutator kernel
is the next open-cover obstruction.
Proposition~\ref{prop:ghost-all-passing-genjac} below
solves these norm equations for every original family
passing the compact elliptic support tests.
